\section{Un coeficiente de orden dos del error de torres}
\paragraph{Teorema revisado, con las entradas uniformes admitidas de las rondas 1, 2 y 8.}
Para las máscaras difusas del apartado 2 y los menores puros de una fila, se tiene
\[
 R(A)=g_T(A)-u_T(A)=k^{-2}F_2(z)+O(k^{-3}),\quad z=kD,
 \qquad |F_2(z)-F_2(z_{\rm hijo})|\le C/k .
\]
Aquí $D$ es el escalado bistocástico exacto de la máscara y las constantes sólo
dependen de los parámetros de difusión y brecha. La prueba siguiente da $F_2$ por extracción
de coeficiente. Trabajar con $z$ usa información del propio punto $P$, pues $D$
se obtiene escalando $P\circ(1-P)$. Por ronda 8, parte 02, $D=p(1+O(k^{-1}))$
entrada a entrada en el soporte y conserva las dos cotas uniformes, con margen menor.
\paragraph{Interpolación en el tamaño.}
En cada lado usamos la medida de probabilidad uniforme sobre sus $k$ índices.
Los productos de núcleos significan integración sobre su índice común. El núcleo $z$
tiene dos márgenes uno y $\|z-1\|_{\op}\le r'<1$. Para un parámetro complejo $s$ escribimos
\[
 v(1-sv)=z\exp(\alpha+\beta),\qquad \int_C v=1,\quad\int_R v=1 .
\]
La ecuación es por entradas, con $v=0$ en los ceros de $z$. En $s=0$ se tiene
$v=z$, $\alpha=\beta=0$. Fijamos $\int_R\alpha=\int_C\beta$. El jacobiano de márgenes
allí tiene bloques diagonales identidad y bloques cruzados $z,z^\top$.
Su inverso en la gauge tiene norma $L^2$ acotada por la brecha. Las ecuaciones
de cada coordenada y $\sup|z|\le C$ dan norma $L^\infty$ acotada del inverso.
Para $|s|<s_0$ fijo, la raíz de $v(1-sv)=w$ que parte de $w$ y el exponencial
tienen derivadas acotadas uniformemente en un entorno fijo. El método del inverso
congelado da allí contracción en los potenciales y una solución holomorfa uniformemente acotada.
En $s=1/k$, $P=sv$ por la estacionariedad, la rama pequeña y la unicidad de torres.
El mismo sistema da esa rama en cualquier replicación de la máscara.
\paragraph{Determinante del operador de líneas de torres.}
Pongamos $\ell=v(1-sv)/(1-2sv)$, $a=\int_C\ell$, $b=\int_R\ell$ y $\bar\ell=\int\ell$.
El núcleo normalizado con la componente constante retirada es
\[
 W=\ell/\sqrt{ab}-\sqrt{ab}/\bar\ell .
\]
Para $s=1/k$ el operador de líneas $S$ tiene grados $a,b$ y conductancias $s\ell$.
Sus bloques diagonales se normalizan a identidad; el bloque cruzado se descompone en
una singularidad uno, con vectores proporcionales a $\sqrt a,\sqrt b$, y el resto $W$.
El pseudodeterminante del normalizado es $2\det(I-W^\top W)$.
En una congruencia $D_0^{1/2} S_{\rm norm}D_0^{1/2}$ de corango uno, los cofactores
dan el factor $\det D_0\,v_0^\top D_0^{-1}v_0$, donde $v_0$ es el nulo unitario previo.
Aquí ese producto cuadrático es $1/\bar\ell$. Luego en el tamaño físico,
\[
 \log\pd S=(\langle\log a\rangle+\langle\log b\rangle)/s+
 \log2+\log\det(I-W^\top W)-\log\bar\ell .
\]
Los corchetes denotan promedios sobre el lado o el producto que corresponda.
Toda transpuesta en prolongaciones complejas es bilineal.
